% !TeX program = pdflatex
% !BIB program = bibtex

\documentclass[a4paper,12pt]{scrreprt}

% --------------------------------------------------------------
% Pakete & Grundeinstellungen
% --------------------------------------------------------------
\usepackage[T1]{fontenc}
\usepackage[utf8]{inputenc}
\usepackage[ngerman]{babel}
\usepackage{setspace}
\usepackage{enumitem}
\usepackage{graphicx}
\usepackage{float}
\usepackage{fancyhdr}
\usepackage{longtable}
\usepackage{booktabs}
\usepackage{hyperref}
\usepackage[absolute,overlay]{textpos}
\usepackage{xcolor,colortbl}
\usepackage{amsmath}
\usepackage[round,authoryear]{natbib}
\usepackage{acronym}
\usepackage{caption}
\usepackage{tikz}
\usepackage{geometry}

% --------------------------------------------------------------
% Layout
% --------------------------------------------------------------
\geometry{left=2.5cm,right=2.5cm,top=2.5cm,bottom=2.5cm}
\setstretch{1}
\sloppy
\setlength{\headheight}{14pt}
\setlength{\parindent}{0pt}
\setlength{\parskip}{1em}
\pagestyle{fancy}
\fancyhf{}
\fancyhead[L]{\textbf{Bachelorarbeit}\\Angewandte Informatik}
\fancyhead[R]{Leon Nöth \\861572}
\fancyfoot[C]{\thepage}

% Textpositionierung
\setlength{\TPHorizModule}{1mm}
\setlength{\TPVertModule}{1mm}

% TikZ Bibliothek
\usetikzlibrary{mindmap,trees}


% Hyperlinkfarben
\hypersetup{
    colorlinks=true,
    linkcolor=black,
    citecolor=blue,
    urlcolor=blue
}

% --------------------------------------------------------------
\begin{document}

\begin{titlepage}
    \thispagestyle{empty} % kein Header/Footer
    
    % Kopfbereich: links Anschrift, rechts Logo
    \begin{textblock*}{7cm}(2cm,3cm)
        \raggedright
        Wilhelm Büchner Hochschule \\
        Fachbereich Medieninformatik \\
        Hilpertstraße 31 \\
        64295 Darmstadt
    \end{textblock*}
    \hfill
    \begin{minipage}[t]{0.3\textwidth}
        \raggedleft
        \includegraphics[width= 5.4cm]{logo.png} 
    \end{minipage}
    
    \vspace*{3cm}
    
    % Titel & Untertitel zentriert
    \centering
    {\LARGE\textbf{Bachelorarbeit}}\\[1cm]
    {\Large Fachbereich Medieninformatik}\\[2cm]
    
    
    \textbf{Konzeption und prototypische Entwicklung eines LLM-basierten Assistenzsystems zur
    Generierung von BPMN-Modellen auf Basis dialogischer Prozesslogik}\\
    
    \vfill
    
    % Ab hier linksbündig
    \raggedright
    \textbf{Leon Nöth} \\ 
    Matrikelnummer: 861572 \\[1cm]
    
    Angewandte Informatik \\[0.3cm]
    Betreut durch: Prof. Dr. Thomas Freytag \\[2cm]
    Abgabetermin: 01.01.2026

    \raggedleft
    \today
    
\end{titlepage}

\newpage
\section*{Gliederung}

\begin{enumerate}[label=\arabic*.]
    \item Einleitung
    \begin{enumerate}[label*=\arabic*.]
        \item Motivation
        \item Zielsetzung und Forschungsfragen
        \item Methodik
        \item Aufbau der Arbeit
    \end{enumerate}
    
    \item Theoretische Grundlagen
    \begin{enumerate}[label*=\arabic*.]
        \item Business Process Model and Notation (BPMN)
            \subitem {2.1.1 Notation und Symbolik}
            \subitem {2.1.2 BPMN im Kontext der Prozessdigitalisierung}
        \item Large Language Models (LLMs) und Prompt-Engineering
            \subitem {2.2.1 Funktionsweise und Potenziale von LLMs}
            \subitem {2.2.2 Prompt Engineering als Steuerungsinstrument}
            \subitem {2.2.3 LLMs im Geschäftsprozesskontext}
        \item Chatbot-Frameworks und Dialogmodellierung
            \subitem {2.3.1 Konversationelle KI und Vertrauen}
            \subitem {2.3.2 Iterative Prozessmodellierung durch Dialog}
            \subitem {2.3.3 Rolle des Feedbacks und adaptiver Rückfragen}
            \subitem {2.3.4 Soziale und ethische Aspekte}
            \subitem {2.3.5 Grenzen automatisierter Dialogsysteme}
    \end{enumerate}
    
    \item Konzepterstellung zur BPMN-Generierung aus Konversation
    \begin{enumerate}[label*=\arabic*]
        \item Ausgangslage und Problemrahmen
        \item Anforderungen an das Assistenzsystem
        \item Extraktion von Prozessdaten aus der Konversation
        \item Mapping-Strategien von Dialogelementen auf BPMN-Elemente
        \item Konzeptionelle Systemarchitektur
    \end{enumerate}
    
	\item Prototypische Implementierung
	\begin{enumerate}[label*=\arabic*.]
		
		\item Technologiestack und Systemumgebung
		\subitem{4.1.1 Godot Engine und GDScript}
		\subitem{4.1.2 LLM-Ausführung mit Ollama}
		\subitem{4.1.3 Modellarchitektur des verwendeten Sprachmodells}
		
		\item Systemarchitektur und Verarbeitungsprozess
		\subitem{4.2.1 Gesamtablauf der Datenverarbeitung}
		\subitem{4.2.2 Userinput → LLM → JSON → Modell}
		\subitem{4.2.3 Übersicht der zentralen Module (ChatController, Loader, NodeFactory)}
		
		\item Dialogverarbeitung und JSON-Generierung
		\subitem{4.3.1 Benutzerinteraktion und Dialogsteuerung}
		\subitem{4.3.2 Prompt-Design für strukturierte Prozessbeschreibungen}
		\subitem{4.3.3 Extraktion von Prozessdaten und JSON-Erzeugung}
		
		\item Erzeugung und Aufbau der BPMN-Modelle
		\subitem{4.4.1 Parsing der JSON-Daten (BpmnJsonLoader)}
		\subitem{4.4.2 Instanziierung der Modellelemente (NodeFactory)}
		\subitem{4.4.3 Verwaltung von Elementbeziehungen und IDs}
		\subitem{4.4.4 Erweiterbarkeit für weitere BPMN-Konstrukte}
		
		\item Layout, Verbindungserzeugung und visuelle Darstellung
		\subitem{4.5.1 Automatische Platzierungsstrategien (LayoutEngine, LayoutHandler)}
		\subitem{4.5.2 Routing der Sequenzflüsse (FlowConnectionHandler, AutoConnector)}
		\subitem{4.5.3 Herausforderungen bei komplexen Prozessmodellen}
		
		\item Fallbeispiele und Generierungsergebnisse
		\subitem{4.6.1 Linearer Prozess}
		\subitem{4.6.2 Entscheidungsstruktur}
		\subitem{4.6.3 Parallele Abläufe}
		
	\end{enumerate}
    
    \item Evaluation
    \begin{enumerate}[label*=\arabic*]
        \item Evaluationsdesign (Qualitätskriterien für BPMN)
        \item Durchführung und Bewertung der Modelle
        \item Reflexion und kritische Analyse
    \end{enumerate}
    
    \item Fazit und Ausblick
    \begin{enumerate}[label*=\arabic*]
        \item Zusammenfassung der Ergebnisse
        \item Beantwortung der Forschungsfragen
        \item Potenziale für Weiterentwicklung und Praxis
    \end{enumerate}
\end{enumerate}

\newpage

\section*{1. Einleitung}

Die jüngsten fortschritte im Bereich der künstlichen Intelligenz (KI) und insbesondere von Large Language Models (LLMs) eröffnen neue Möglichkeiten zur automatisierten Verarbeitung natürlicher Sprache. Im Geschäftsprozessmanagement besteht dabei besonderes Potenzial, da die Erstellung von BPMN-Modellen trotz ihres etablierten Standards weiterhin zeitintensiv ist und spezialisiertes Fachwissen erfordert.

Diese Arbeit untersucht die Konzeption und prototypische Umsetzung eines Assistenzsystems, das mithilfe eines LLMs aus dialogischen Prozessbeschreibungen automatisch BPMN-Modelle generiert. Ziel ist es, den Übergang von informellen Beschreibungen zu formalisierten Prozessmodellen zu unterstützen.

\section*{1.1 Motivation}

Die rasante Entwicklung von Large Language Models (LLMs) und künstlicher Intelligenz (KI) stellt einen der spannendsten technologischen Fortschritte unserer Zeit dar. Schon seit Beginn meines Studiums habe ich ein besonderes Interesse an den Möglichkeiten dieser Technologien entwickelt und kontinuierlich in mein Wissen sowie praktische Erfahrungen in diesem Bereich investiert.

Die Potenziale, die moderne KI-Systeme insbesondere im Kontext der Automatisierung, Optimierung und Unterstützung von Geschäftsprozessen bieten, faszinieren mich dabei besonders.

Die Verbindung von LLMs mit BPMN-Prozessen erscheint mir besonders reizvoll, da ich bereits umfassende praktische Erfahrungen in wirtschaftlichen und industriellen Bereichen sammeln konnte. Während meiner Tätigkeiten in verschiedenen Branchen habe ich gesehen, dass nahezu jedes Unternehmen auf strukturierte Prozesse angewiesen ist, um effizient und effektiv arbeiten zu können.

Prozesse klar zu definieren, zu modellieren und zu dokumentieren hat sich stets als entscheidend für den Erfolg und die Qualität der Arbeit erwiesen. Dabei konnte ich nicht nur die theoretische Relevanz von Prozessmanagement erkennen, sondern auch unmittelbar erleben, wie gut definierte Abläufe die Produktivität steigern und die Kommunikation innerhalb von Teams verbessern.

Die Idee, ein KI-gestütztes System zu entwickeln, das auf Basis natürlicher Sprache automatisch Prozessmodelle generiert, vereint meine technische Begeisterung für KI mit meinem Interesse an der Optimierung von Arbeitsabläufen. Sie bietet die Möglichkeit, komplexe betriebliche Zusammenhänge zugänglicher zu machen und die Prozessmodellierung zu vereinfachen. 

Ziel dieses Projektes ist es daher, die Potenziale moderner KI-Systeme gezielt auf die Prozessmodellierung zu übertragen und dabei sowohl die Effizienz der Modellierung als auch die Qualität der erzeugten BPMN-Darstellungen zu steigern.

Darüber hinaus motiviert mich die Aussicht, in einem aktuellen und zukunftsrelevanten Forschungsfeld tätig zu sein. Die kontinuierliche Weiterentwicklung von LLMs und KI-Anwendungen eröffnet zahlreiche neue Anwendungsfelder, und ich möchte aktiv daran mitarbeiten, diese Technologien auf praxisrelevante Fragestellungen anzuwenden.

Dieses Projekt bietet mir die Gelegenheit, sowohl meine fachlichen Kompetenzen zu erweitern als auch einen Beitrag zu einem zukunftsweisenden Themenbereich zu leisten, der in den kommenden Jahren voraussichtlich eine zentrale Rolle in Wirtschaft und Gesellschaft spielen wird.

\section*{1.2 Zielsetzung und Forschungsfragen}

Ziel dieser Bachelorarbeit ist die Konzeption und prototypische Entwicklung eines Assistenzsystems, das mithilfe eines Large Language Models (LLM) in der Lage ist, textbasierte Prozessbeschreibungen automatisch in Business Process Model and Notation (BPMN)-Diagramme zu überführen. Auf diese Weise soll untersucht werden, inwiefern aktuelle LLM-Technologien geeignet sind, komplexe sprachliche Eingaben so zu strukturieren, dass daraus formal gültige und semantisch sinnvolle Prozessmodelle entstehen. 

Im Mittelpunkt steht dabei die Entwicklung eines Konzepts, das die Fähigkeiten moderner Sprachmodelle mit den Anforderungen der Prozessmodellierung verbindet. Durch die prototypische Umsetzung soll überprüft werden, welche Ansätze, Architekturen und Prompt-Strategien sich besonders eignen, um eine möglichst präzise und konsistente Modellgenerierung zu ermöglichen. Die Arbeit verfolgt somit sowohl ein theoretisches als auch ein praktisches Ziel: Einerseits sollen die Potenziale und Grenzen von LLMs im Kontext der BPMN-Modellierung analysiert werden, andererseits soll eine funktionsfähige technische Lösung als Machbarkeitsnachweis entstehen.

\newpage

Aus dieser Zielsetzung ergeben sich die folgenden Forschungsfragen \textit{Fn}:

\begin{itemize}
    \item[\textit{F1}\textbf{:}] Wie kann ein Large Language Model so eingesetzt werden, dass es textbasierte Prozessbeschreibungen zuverlässig in strukturierte BPMN-Modelle überführt?
    \item[\textit{F2}\textbf{:}] Welche Methoden, Prompt-Strategien und architektonischen Ansätze sind geeignet, um die Genauigkeit und Qualität der automatischen Prozessmodellierung zu verbessern?
    \item[\textit{F3}\textbf{:}] Welche Chancen und Grenzen ergeben sich aus dem Einsatz von LLMs für die Modellierung von Geschäftsprozessen im Vergleich zu traditionellen Vorgehensweisen?
\end{itemize}

Die Beantwortung dieser Fragen soll Aufschluss darüber geben, in welchem Umfang LLMs bereits heute einen Mehrwert für die Prozessmodellierung bieten können und welche Entwicklungen zukünftig erforderlich sind, um den praktischen Einsatz solcher Systeme in betrieblichen Kontexten zu ermöglichen.

\section*{1.3 Methodik}

Die vorliegende Arbeit verfolgt einen methodischen Ansatz, der sowohl theoretische als auch praktische Elemente umfasst. Ziel ist es, durch eine fundierte Literaturrecherche zunächst den aktuellen Forschungsstand zu erfassen, um darauf aufbauend ein eigenes Konzept sowie einen funktionsfähigen Prototypen zu entwickeln. Die Methodik orientiert sich an den Prinzipien wissenschaftlicher Arbeit im Bereich der angewandten Informatik und kombiniert Analyse mit experimenteller Umsetzung.

\subsection*{1.3.1 Literaturrecherche}
Die Erarbeitung der theoretischen Grundlagen erfolgte primär durch eine systematische Literaturrecherche. Hierbei wurden wissenschaftliche Publikationen, Konferenzbeiträge und Fachartikel aus relevanten Datenbanken und digitalen Bibliotheken ausgewertet. Besonders genutzt wurden hierfür \textit{Google Scholar}, die \textit{SpringerLink}-Bibliothek sowie die \textit{EBSCOhost Research Database}. Die Suche erfolgte anhand gezielter Schlagwörter wie \textit{"Large Language Models"}, \textit{"BPMN"}, \textit{"Process Modeling"}, \textit{"Prompt Engineering"} und \textit{"LLM Process Modeling"}. Ergänzend wurde Fachliteratur herangezogen, um bestehende Konzepte und technologische Entwicklungen im Bereich der Prozessmodellierung und künstlichen Intelligenz zu betrachten,analysieren und einzuordnen.

\newpage

\subsection*{1.3.2 Konzeption und prototypische Umsetzung}
Auf Grundlage der analysierten Literatur wurde ein konzeptioneller Ansatz entwickelt, der die automatische Generierung von BPMN-Modellen durch ein Large Language Model ermöglicht. Anstelle einer herkömmlichen Framework-Integration erfolgt die Umsetzung in Form einer eigenständigen Anwendung, die mit der Game Engine Godot 4.5.1 realisiert wurde. Das System umfasst ein integriertes Dialoginterface, über das Nutzerinnen und Nutzer mithilfe eines selbstgehosteten LLMs interaktiv Prozessbeschreibungen formulieren können.

Das LLM erzeugt auf Basis eines definierten Masterprompts eine strukturierte JSON-Repräsentation des gewünschten Prozessmodells. Diese wird anschließend von mehreren speziell entwickelten Komponenten, darunter ein LayoutHandler, eine LayoutEngine, ein FlowHandler sowie ein BPMNJsonLoader, verarbeitet und in ein visuell darstellbares BPMN-Diagramm überführt. Der entwickelte Prototyp demonstriert damit eine vollständige Systemlösung, die den dialogbasierten Übergang von natürlicher Sprache zu formalisierten BPMN-Modellen innerhalb einer konsistenten Anwendung realisiert.

\subsection*{1.3.3 Evaluation}
Zur Bewertung des entwickelten Prototyps ist eine Analyse vorgesehen. Dabei wird untersucht, inwiefern die generierten BPMN-Modelle den definierten Anforderungen an Struktur, Vollständigkeit und semantische Korrektheit entsprechen. Die Ergebnisse sollen im Hinblick auf die eigenen formulierten Forschungsfragen interpretiert und kritisch reflektiert werden.

Durch diese Punkte stellt die Methodik sicher, dass sowohl der theoretische als auch der praktische Teil der Arbeit wissenschaftlich fundiert und nachvollziehbar durchgeführt wird. Sie erlaubt eine strukturierte Herangehensweise an das Forschungsproblem und gewährleistet, dass die gewonnenen Erkenntnisse transparent, reproduzierbar und methodisch begründet sind.

\section*{1.4 Aufbau der Arbeit}

Die vorliegende Arbeit ist in sechs Hauptkapitel gegliedert.

Kapitel~1 führt in das Thema ein, erläutert die Motivation, formuliert die Zielsetzung sowie die zugrunde liegenden Forschungsfragen und beschreibt die verwendete Methodik.

In Kapitel~2 werden die theoretischen Grundlagen dargestellt. Dazu zählen die Business Process Model and Notation (BPMN), die Funktionsweise von Large Language Models (LLMs) einschließlich relevanter Aspekte des Prompt-Engineerings sowie die Grundlagen der Chatbot-Frameworks und der Dialogmodellierung.

Kapitel~3 beschreibt das konzeptionelle Vorgehen zur automatischen Generierung von BPMN-Diagrammen auf Basis dialogischer Prozesslogik. Hier werden die Anforderungen an das Assistenzsystem, die Extraktion strukturierter Prozessdaten aus Konversationen, die Abbildung auf BPMN-Elemente sowie die Systemarchitektur erläutert.

Kapitel~4 widmet sich der prototypischen Implementierung des entwickelten Konzepts. Es behandelt den eingesetzten Technologiestack, die Umsetzung des BPMN-Generators, die Integration in ein dialogbasiertes Assistenzsystem sowie beispielhafte Prozessgenerierungen.

In Kapitel~5 erfolgt die Evaluation des entwickelten Prototyps. Dabei werden das Evaluationsdesign und die Qualitätskriterien für BPMN-Modelle beschrieben, die Ergebnisse ausgewertet und kritisch reflektiert.

Abschließend fasst Kapitel~6 die wesentlichen Erkenntnisse zusammen, beantwortet die Forschungsfragen und gibt einen Ausblick auf mögliche Weiterentwicklungen und praktische Einsatzpotenziale.

\section*{2. Theoretische Grundlagen}

Dieses Kapitel erläutert die theoretischen Grundlagen, auf denen das in dieser Arbeit entwickelte 
LLM-basierte Assistenzsystem zur automatischen Generierung von BPMN-Modellen aufbaut. 
Ziel ist es, die relevanten Konzepte aus den Bereichen Prozessmodellierung, Sprachmodelltechnologie 
und dialogorientierte Interaktionssysteme systematisch darzustellen. 

Dementsprechend ist das Kapitel in drei Kernbereiche gegliedert:
\begin{enumerate}
	\item \ac{BPMN} als Modellierungsstandard,
	\item \acp{LLM} und Prompt-Engineering als technologische Basis sowie
	\item Chatbot-Frameworks und Dialogmodellierung als interaktives Vermittlungsprinzip.
\end{enumerate}

\section*{2.1 Business Process Model and Notation (BPMN)}
\subsubsection*{2.1.1 Notation und Symbolik}

Die \ac{BPMN} stellt einen international anerkannten Standard zur Modellierung, Visualisierung 
und Analyse von Geschäftsprozessen dar. Ursprünglich wurde sie 2002 bei IBM entwickelt 
und durch die Business Process Management Initiative (BPMI) veröffentlicht. 
Im Jahr 2005 übernahm die Object Management Group (OMG) die BPMI und führte BPMN 
als offiziellen OMG-Standard weiter, wodurch die Notation in die Reihe etablierter 
Modellierungssprachen wie der \ac{UML} eingegliedert wurde 
\citep[vgl.][]{ElstermannEtAl2023,OMG2013}. 

Die Symbolik der \ac{BPMN} basiert auf klar definierten Flussobjekten: 
Ereignisse werden als Kreise, Aktivitäten als abgerundete Rechtecke und Gateways als Rauten dargestellt. 
Diese Elemente sind über Sequenzflüsse miteinander verbunden, welche den Ablauf und die 
Abhängigkeiten zwischen Prozessschritten visualisieren 
\citep[vgl.][S.~85--87]{Dumas2021}. 
Gateways dienen dabei der Steuerung von Verzweigungen und Zusammenführungen. 
Unterschieden wird zwischen exklusiven (XOR-), parallelen (AND-) und inklusiven (OR-)Gateways 
\citep[vgl.][S.~90--95]{ElstermannEtAl2023}. 

\begin{figure}[H]
	\centering
	\includegraphics[width=\textwidth]{BPMN Allgemein vers.1.0.1.png}
	\caption{Grundelemente der BPMN-Notation (eigene Darstellung in Anlehnung an \citealp[S.~85~ff.]{ElstermannEtAl2023})}
	\label{fig:bpmn_notation}
\end{figure}

Ein wesentlicher Vorteil der \ac{BPMN} liegt in ihrer semantischen Tiefe: 
Sie ermöglicht nicht nur die grafische Darstellung, sondern auch die Transformation in 
maschinenlesbare Prozesssprachen wie \ac{BPEL}, wodurch eine automatisierte Ausführung 
modellierter Prozesse möglich wird \citep[vgl.][S.~2--4]{mazanek2011}. 
Wie \citet[S.~22--27]{ElstermannEtAl2023} hervorheben, sollte die Modellierung dabei stets 
so gestaltet sein, dass sie über eine reine Beschreibung des Ablaufs hinausgeht und eine 
fundierte Analyse und spätere Verbesserung des Prozesses ermöglicht. 
Der Detaillierungsgrad eines Prozessmodells hängt vom jeweiligen Verwendungszweck ab. 
Während einfache Modelle der Orientierung und Kommunikation dienen, erfordern 
analytische oder automatisierungsorientierte Modelle zusätzliche Informationen, 
etwa zu Durchlaufzeiten, Eingabe- und Ausgabeformaten oder potenziellen Fehlerquellen. 
Damit wird die Prozessmodellierung zum integralen Bestandteil des gesamten 
BPM-Lebenszyklus, der von der Identifikation und Analyse über die Verbesserung 
bis hin zur Implementierung und Überwachung reicht. 
Die Qualität und Tiefe des Modells entscheidet somit maßgeblich darüber, 
inwieweit ein Prozess automatisiert, verbessert oder digital unterstützt werden kann, 
ohne dabei technische Implementierungsdetails vorwegzunehmen 
\citep[vgl.][S.~22--27]{ElstermannEtAl2023}.


\subsubsection*{2.1.2 BPMN im Kontext der Prozessdigitalisierung}

Im Zuge der digitalen Transformation hat sich die \ac{BPMN} in der industriellen Praxis als vielseitiges Instrument etabliert, das von der reinen Dokumentation bis hin zur (teil-)automatisierten Ausführung von Geschäftsprozessen reicht \citep[vgl.][Kap.~3.5.8]{ElstermannEtAl2023}. 
Damit eignet sich die Notation, Prozesse konsistent, verständlich und bei Bedarf maschinenlesbar abzubilden.

Auf konzeptioneller Ebene unterstützt die Modellierung zudem die Reduktion organisatorischer und technischer Schnittstellen. Weniger Schnittstellen verringern Medienbrüche, Doppelarbeit und Fehler sowie Zeitverluste und Kosten \citep[vgl.][S.~226]{ElstermannEtAl2023}. 
Die \ac{BPMN} trägt hier zur Transparenz über Aufgaben, Zuständigkeiten und Informationsflüsse bei.

Bereits in der Modellbildung sollten Aspekte berücksichtigt werden, die eine spätere digitale Ausführung erleichtern, allerdings ohne Implementierungsdetails vorwegzunehmen. Je präziser die Prozesse beschrieben sind, desto leichter fällt die Umsetzung \citep[vgl.][S.~10]{ElstermannEtAl2023}.

Mit dem Aufkommen generativer \ac{KI} erfährt die Prozessmodellierung einen grundlegenden Wandel. 
\citet[S.~81--85]{KampikWarmuthRebmannLeiteKirchner2025} beschreiben, dass Prozessdarstellungen zunehmend von statischen, beschreibenden Modellen zu dynamischen, lernfähigen Wissensrepräsentationen übergehen. 
Zentral ist dabei das Konzept der Large Process Models (LPMs), welche die statistische Ausdrucksstärke von \acp{LLM} mit der logischen Struktur wissensbasierter Systeme verbinden. 
Dadurch entsteht ein neues Verständnis von Prozessmanagement, bei dem Modelle nicht mehr ausschließlich der statischen Abbildung bestehender Abläufe dienen, sondern potenziell zur Grundlage adaptiver und kontextsensitiver Analysen werden. 
In Anlehnung an \citet{KampikWarmuthRebmannLeiteKirchner2025} lässt sich daraus schließen, dass zukünftige Prozesssysteme zunehmend fähig sein könnten, auf Prozessdaten zu reagieren und kontextbasierte Handlungsempfehlungen zu unterstützen.

In Anlehnung an \citet{ElstermannEtAl2023} und \citet{KampikWarmuthRebmannLeiteKirchner2025} lässt sich daraus ableiten, dass die \ac{BPMN} im Kontext der Prozessdigitalisierung nicht nur als Visualisierungsstandard, sondern als methodisches Rückgrat einer wissens- und datengetriebenen Prozessarchitektur fungiert, in der klassische Modellierung zunehmend mit \ac{KI}-gestützter Analyse und Unterstützung verschmilzt. 
Damit verschiebt sich der Fokus von der statischen Modellierung hin zu einer kontinuierlich lernenden Prozessarchitektur, die neue Anforderungen des digitalen Zeitalters integriert.

\section*{2.2 Large Language Models (LLMs) und Prompt-Engineering}

\subsubsection*{2.2.1 Funktionsweise und Potenziale von LLMs}

\acp{LLM} sind neuronale Sprachmodelle, die auf der Verarbeitung und Generierung natürlicher Sprache beruhen. 
Sie basieren auf sogenannten Transformer-Architekturen, die durch Selbstaufmerksamkeitsmechanismen (self-attention) in der Lage sind, semantische Abhängigkeiten über lange Textsequenzen hinweg zu erfassen. 
Dadurch können sie sprachliche Muster, Syntax und Bedeutungszusammenhänge präzise modellieren und kontextabhängige Antworten generieren. 
\citet{OpenAI2023} beschreiben, dass GPT-4 über mehrere hundert Milliarden Parameter verfügt und mithilfe von Reinforcement Learning from Human Feedback (RLHF) so trainiert wird, dass es logisch konsistente, kontextbewusste und sprachlich kohärente Ergebnisse liefert. 
Ebenso zeigen \citet{Touvron2023}, dass LLaMA-2-Modelle in verschiedenen Größen bis zu 70 Milliarden Parametern skalieren und eine besonders hohe Effizienz bei der Wissensrepräsentation natürlicher Sprache aufweisen.

Die Stärke dieser Modelle liegt in ihrer Fähigkeit, aus offenen Textkorpora latente Wissensstrukturen zu extrahieren und kontextbezogene Antworten zu generieren. 
Dadurch können sie nicht nur sprachliche, sondern auch semantische und logische Beziehungen zwischen Konzepten erkennen und wiedergeben. 
Aktuelle Forschungsarbeiten wie \citep[S.~99--100]{Moeller2025} zeigen, dass \acp{LLM} zunehmend zur Analyse und Modellierung von Geschäftsprozessen eingesetzt werden können. 
So gelingt es, Ereignisstrukturen und Prozesslogiken aus unstrukturierten Daten zu extrahieren und in formale Modelle zu überführen. 
Dies eröffnet neue Potenziale für die automatisierte Prozessanalyse, Prozessdokumentation und Entscheidungsunterstützung.

Gleichzeitig verweisen die Autoren auf bestehende Grenzen: 
Generische Sprachmodelle verfügen oftmals nicht über domänenspezifisches Wissen und neigen zu semantischen Fehlern, wenn die Trainingsdaten keine prozessrelevanten Muster enthalten. 
Die Qualität der Ergebnisse hängt daher stark von der Datenbasis, der Prompt-Gestaltung und der Kontextsteuerung ab. 
Insgesamt markieren \acp{LLM} damit einen Paradigmenwechsel in der Sprachverarbeitung, weg von regelbasierten Textanalysen hin zu generativen, wissensbasierten Systemen, die Sprache als dynamische Wissensquelle interpretieren und aktiv zur Strukturierung organisationaler Wissensprozesse beitragen.

\subsubsection*{2.2.2 Prompt Engineering als Steuerungsinstrument}

Prompt Engineering beschreibt die systematische Gestaltung und Strukturierung von Eingaben an ein Sprachmodell, um dessen Ausgabequalität gezielt zu steuern. 
\citet[S.~485--489]{Hsu2024} definieren es als bewussten Entwurfsprozess sprachlicher Instruktionen, der Parameter wie Aufgabenbeschreibung, Formatvorgaben und Kontextinformationen kombiniert, um ein gewünschtes Antwortverhalten zu erzielen. 
Zu den zentralen Techniken zählen dabei klare Instruktionen (Instruction Design), die Bereitstellung von Beispielen (Few-Shot Prompting) und die iterative Anpassung über Rückkopplungsschleifen (Refinement) \citep[S.~486--488]{Hsu2024}. 

Im Kontext der prozessbezogenen Modellgenerierung zeigen empirische Studien, dass die Struktur und Präzision von Prompts einen messbaren Einfluss auf die Qualität der erzeugten Modelle besitzen. 
\citet[S.~5--7]{Kourani2024} heben hervor, dass Strategien wie Role Prompting, Knowledge Injection und Negative Prompting die semantische Konsistenz der generierten BPMN-Modelle erhöhen. 
Ebenso zeigen \citet[S.~3--4]{KoepkeSafan2024}, dass standardisierte Promptvorlagen und JSON-Zwischenrepräsentationen die Modellqualität verbessern und die Rechenlast reduzieren. 

Für diese Arbeit bedeutet dies, dass Prompts präzise Rollen- und Aufgabenhinweise, definierte Ausgabestrukturen (z.\,B. BPMN-XML) und Validierungskriterien enthalten sollten. 
Iteratives Nachschärfen durch gezielte Rückfragen zu Gateways oder Sequenzen ermöglicht eine schrittweise Qualitätssteigerung. 
Prompt Engineering fungiert somit als methodisches Steuerungsinstrument, das generative Sprachmodelle aktiv in den Dienst strukturierter Prozessmodellierung stellt.


\subsubsection*{2.2.3 LLMs im Geschäftsprozesskontext}

Im Geschäftsprozessmanagement übernehmen \acp{LLM} zunehmend die Rolle semantischer Schnittstellen. 
Sie können unstrukturierte Eingaben, etwa Konversationen oder Dokumente, in strukturierte Repräsentationen überführen 
und damit Prozesswissen zugänglich machen \citep[S.~2--3]{Kourani2024}. 
Gleichzeitig zeigen aktuelle Studien, dass rein statistische Modelle ohne geeignete Leitplanken 
zu Inkonsistenzen und mangelnder Erklärbarkeit neigen. 
Hybride Architekturen, in denen \acp{LLM} mit formalem Prozesswissen oder wissensbasierten Komponenten gekoppelt werden, adressieren dieses Problem. 
Ein Beispiel sind die Large Process Models (LPMs), die die probabilistische Sprachkompetenz generativer Modelle 
mit der logischen Präzision wissensbasierter Systeme verbinden 
und dadurch nachvollziehbarere sowie robustere Ergebnisse ermöglichen \citep[S.~81--84]{KampikWarmuthRebmannLeiteKirchner2025}. 

Empirische Untersuchungen zeigen zudem, dass leistungsfähigere Modelle 
komplexe Abhängigkeiten besser erfassen und konsistentere Strukturen erzeugen, 
während kleinere Modelle früh an ihre Grenzen stoßen \citep[S.~6--8]{Kourani2024}. 
Dialogische Frameworks wie \emph{BPMNGen} demonstrieren darüber hinaus, dass Interaktion und gezieltes Prompting 
die syntaktische, semantische und pragmatische Qualität generierter BPMN-Modelle verbessern können. 
Das Framework ermöglicht es, Prozessmodelle iterativ über natürliche Sprache zu generieren und zu verfeinern, 
wobei Qualität über strukturierte Eingaben, JSON-Zwischenrepräsentationen und Rückkopplungsschleifen gesteuert wird 
\citep[S.~27--31]{Hoerner2025}. 

Damit verdeutlicht sich, dass \acp{LLM} im Geschäftsprozesskontext 
nicht nur Werkzeuge der Automatisierung darstellen, sondern zunehmend integrale Bestandteile einer wissensbasierten Prozessarchitektur werden. 
Ihr Einsatz markiert den Übergang von statischer Modellbeschreibung hin zu einer interaktiven, lernfähigen und erklärbaren Prozessintelligenz, 
die den Menschen nicht ersetzt, sondern in seiner Analyse- und Gestaltungsrolle stärkt.


\section*{2.3 Chatbot-Frameworks und Dialogmodellierung}

\subsubsection*{2.3.1 Konversationelle KI und Vertrauenswürdigkeit}

Konversationssysteme bilden eine zentrale Schnittstelle zwischen Mensch und Maschine, insbesondere in der interaktiven Gestaltung digitaler Geschäftsprozesse. 
Dabei umfasst ihr Potenzial nicht nur funktionale, sondern auch soziale Dimensionen, die das Nutzungserlebnis und die Akzeptanz maßgeblich prägen. 
\citet[S.~3--5]{FigueroaTorres2024} unterscheidet drei gesellschaftliche Dimensionen von Chatbots: 
ihre wissenschaftliche Erforschung, ihre wirtschaftliche Nutzung und ihre soziale beziehungsweise persönliche Wirkung. 
Vertrauen spielt dabei in allen drei Bereichen eine zentrale Rolle, insbesondere im Hinblick auf Transparenz und Nachvollziehbarkeit der Systementscheidungen. 

Auch im Kontext digitaler Prozesssysteme gilt Transparenz als Voraussetzung für Vertrauen. 
\citet[S.~220--225]{ElstermannEtAl2023} betonen im Rahmen ihres Konzepts der wertegeleiteten Interaktion, 
dass technologische Systeme nur dann Akzeptanz erzeugen, wenn ihre Funktionsweise und Grenzen für die Beteiligten nachvollziehbar bleiben. 
Eine klare Kommunikation von Annahmen und Vorschlägen, etwa in dialoggestützten Prozessmodellen, 
unterstützt somit eine partizipative und wertorientierte Nutzung.

\subsubsection*{2.3.2 Iterative Modellierung durch natürliche Sprache}

Ein wesentlicher Vorteil konversationeller Prozessmodellierung liegt in der Möglichkeit, Modelle iterativ und kontextsensitiv zu entwickeln. 
\citet[S.~27--31]{Hoerner2025} beschreibt mit \textit{BPMNGen} ein LLM-basiertes Framework, das es erlaubt, 
BPMN-Modelle über natürliche Sprache schrittweise zu generieren und zu verfeinern. 
Durch die dialogische Interaktion können Unklarheiten im Prozessverständnis aufgedeckt und durch gezielte Rückfragen geklärt werden. 
Damit verschiebt sich der Modellierungsprozess von einer rein dokumentarischen zu einer adaptiven, partizipativen Form. 


\subsubsection*{2.3.3 Rolle des Feedbacks und adaptiver Rückfragen}

\citet[S.~425--430]{Klievtsova2023} zeigen in ihrer systematischen Analyse, dass moderne Chatbot-Frameworks potenziell alle Phasen des Business Process Management (\ac{BPM}) unterstützen können, auch wenn derzeit nur Teilaspekte praktisch umsetzbar sind. Durch interaktives Feedback und kontinuierliche Dialoge tragen Chatbots nicht nur zur initialen Modellgenerierung bei, sondern unterstützen auch die schrittweise Optimierung und Qualitätssicherung bestehender Modelle. Die Autoren skizzieren Chatbots als aufkommende Akteure im BPM-Kontext. Denn sie agieren nicht nur als Übersetzer zwischen natürlicher Sprache und formalen Modellen, sondern fungieren zugleich als reflektierende Partner im Qualitäts- und Verbesserungsprozess.

\subsubsection*{2.3.4 Soziale und ethische Implikationen konversationeller Modellierung}

Neben den technischen Potenzialen sind auch soziale und ethische Aspekte der konversationellen Modellierung relevant. 
Systeme, die in geschäftskritischen Entscheidungsprozessen eingesetzt werden, müssen nachvollziehbar agieren und dürfen nicht den Eindruck autonomer Entscheidungsgewalt erwecken. 

\citet[S.~10--13]{FigueroaTorres2024} betont, dass die Zuschreibung von Verantwortung eng mit der wahrgenommenen Transparenz und Kompetenz eines Systems verknüpft ist. 
Nur wenn Nutzer nachvollziehen können, wie Entscheidungen zustande kommen, kann Verantwortung klar zugeordnet werden. 
Diese Forderung lässt sich mit dem Paradigma der wertegeleiteten Interaktion in Beziehung setzen, das \citet[S.~216--226]{ElstermannEtAl2023} als normativen Rahmen formulieren: 
Menschliche Werte wie Verantwortung, Partizipation und Verständlichkeit sollen dabei als leitende Prinzipien für den Einsatz konversationeller Systeme dienen.

\subsubsection*{2.3.5 Grenzen automatisierter Dialogsysteme}

Trotz der beachtlichen Fortschritte moderner Sprachmodelle bleibt der Mensch zentrale Referenz und letztverantwortlicher Akteur im Geschäftsprozessmanagement. 
\citet[S.~4--9]{Kourani2024} zeigen anhand eines LLM-basierten Rahmenwerks zur Prozessmodellierung, dass \acp{LLM} vor allem dann einen Mehrwert bieten, 
wenn sie domänenspezifisches Wissen strukturieren, in kontrollierte Entscheidungsarchitekturen eingebettet sind und über Feedback- sowie Fehlerbehandlungsschleifen in den Modellierungsprozess integriert werden. 

Eine autonome Entscheidungsfindung durch solche Systeme ist hingegen weder intendiert noch verantwortbar: Die Studie weist auf semantische Verständnislücken, Prompt-Sensitivität, Halluzinationen sowie fehlende Erklärbarkeit hin \citep[S.~11--14]{Kourani2024}, 
sodass Ergebnisse ohne menschliche Prüfung und Steuerung nicht verlässlich sind. 
\citet[S.~81--84]{KampikWarmuthRebmannLeiteKirchner2025} skizzieren mit dem Konzept der Large Process Models eine Vision zukünftiger BPM-Systeme, in denen \acp{LLM} mit wissensbasierten Verfahren kombiniert werden, um kontextsensitive und erklärbare Prozessunterstützung zu ermöglichen; die epistemische Kontrolle verbleibt dabei grundsätzlich beim Menschen.

Damit wird deutlich, dass der Mensch auch im Zeitalter generativer KI als leitende Instanz für Analyse, Interpretation und ethische Bewertung erhalten bleiben muss. 
\acp{LLM} fungieren als Werkzeuge, die kognitive und analytische Prozesse ergänzen, nicht jedoch als autonome Entscheider.



\section*{3. Konzepterstellung zur BPMN-Generierung aus Konversation}

Dieses Kapitel entwickelt das konzeptionelle Fundament eines Assistenzsystems, das natürlichsprachliche Dialoge in formal strukturierte BPMN-Modelle überführt. Aufbauend auf den theoretischen Grundlagen aus Kapitel 2 werden die Herausforderungen natürlicher Sprache für die Prozessmodellierung herausgearbeitet, konzeptionelle Anforderungen formuliert und Methoden für die Bedeutungserschließung, Strukturierung und Repräsentation von Prozesswissen entwickelt. Ziel ist es, einen kohärenten Ansatz zu definieren, der die Offenheit dialogischer Interaktion mit den formalen Regeln der Geschäftsprozessmodellierung verbindet und damit die Grundlage für die spätere Implementierung bildet.


\subsection*{3.1 Ausgangslage}

Obwohl in den vergangenen Jahren mehrere Ansätze zur KI-gestützten Prozessmodellierung entstanden sind, zeigt eine Analyse des aktuellen Forschungs- und Entwicklungsstands deutliche funktionale und methodische Lücken. Klassische Werkzeuge wie der Camunda Modeler\footnote{\url{https://camunda.com/platform/modeler/}} sowie Konverter wie der \textit{camunda-bpmn-converter}\footnote{\url{https://github.com/lzgabel/camunda-bpmn-converter}} bieten zwar Möglichkeiten zur Transliteration bestehender Strukturdateien in BPMN-Notation, verfügen jedoch über keine Mechanismen zur semantischen Interpretation natürlicher Sprache. Sie ermöglichen damit keinen Übergang von unstrukturierten Dialogen zu formalen Prozessmodellen, sondern setzen bereits vorstrukturierte Eingabedaten voraus.

Neuere LLM-basierte Ansätze wie der \textit{BPMN Assistant}\footnote{\url{https://github.com/jtlicardo/bpmn-assistant}; vgl. auch \url{https://arxiv.org/abs/2509.24592}} zeigen hingegen, dass eine dialogische Interaktion mit einem Sprachmodell grundsätzlich geeignet ist, Prozessstrukturen automatisch zu erzeugen. Allerdings verdeutlichen sowohl das Projekt als auch die dazugehörige wissenschaftliche Veröffentlichung, dass die derzeit verfügbaren Systeme funktionale Einschränkungen aufweisen. Beispielsweise unterstützt der BPMN Assistant lediglich eine begrenzte Untermenge der BPMN-Elementtypen, verfügt über kein konsistentes Layout für Pools oder Lanes und kann manuelle Veränderungen am Diagramm nicht in seine internen Repräsentationen einbeziehen. Modellierungsentscheidungen erfolgen ausschließlich auf Basis der jeweils letzten vom System generierten Version, was eine inkrementelle, dialogisch abgestimmte Prozessentwicklung erheblich erschwert.

Kommerzielle Lösungen wie der Camunda BPMN Copilot\footnote{\url{https://docs.camunda.io/docs/components/early-access/alpha/bpmn-copilot/}} integrieren Sprachmodelle in professionelle Modellierungsumgebungen, folgen jedoch einer proprietären und cloudbasierten Architektur. Dies begrenzt ihre wissenschaftliche Nutzbarkeit sowohl im Hinblick auf Reproduzierbarkeit als auch im Hinblick auf Datenschutz, konfigurierbare Modellkontexte und die Möglichkeit, alternative lokale Sprachmodelle einzubinden. Darüber hinaus bieten solche Systeme in der Regel keinen direkten Zugriff auf interne Dialog- oder Modellierungsdaten, wodurch eine empirische Analyse von Interaktionsmustern oder Fehlerszenarien nicht möglich ist.

Zusammenfassend zeigt sich, dass existierende Systeme entweder auf formale Konvertierung begrenzt sind, nur Teilmengen der BPMN-Notation unterstützen oder aufgrund proprietärer Beschränkungen nicht für experimentelle und wissenschaftliche Weiterentwicklung geeignet sind. Es fehlt bislang ein offen gestaltetes, lokal ausführbares und modular erweiterbares Assistenzsystem, das dialogische Prozessbeschreibungen vollständig interpretiert, strukturiert und iterativ verfeinert und zugleich die wissenschaftliche Analyse der Interaktion ermöglicht. Diese Lücke bildet die konzeptionelle Ausgangsbasis für die in dieser Arbeit entwickelte Systemarchitektur.


\subsection*{3.2 Anforderungen an das Assistenzsystem}

Das entwickelte Assistenzsystem verfolgt das Ziel, Nutzende bei der Modellierung von Geschäftsprozessen zu unterstützen, indem frei formulierte Spracheingaben in formal definierte BPMN-Strukturen überführt werden. Hieraus ergeben sich sowohl funktionale als auch nicht-funktionale Anforderungen, die im Zusammenspiel bestimmen, welche Fähigkeiten das System aufweisen muss, um den in Kapitel 2 dargestellten Herausforderungen dialogischer Prozessmodellierung angemessen zu begegnen.

Funktional benötigt das System die Fähigkeit, natürliche Sprache so zu interpretieren, dass sie als strukturiertes Prozesswissen nutzbar wird. Dies umfasst nicht nur die Identifikation einzelner Aktivitäten, Ereignisse oder Rollen, sondern insbesondere die Fähigkeit, sprachliche Äußerungen im jeweiligen Gesprächskontext zu verstehen. Da Prozesswissen selten vollständig oder linear beschrieben wird, muss das System iterative und unvollständige Eingaben erkennen und auf dieser Grundlage gezielte Rückfragen formulieren können. Die Generierung von BPMN-Modellen darf daher nicht als eine reine einfache Umwandlung verstanden werden, sondern als dialogischer Prozess, der semantische Lücken sichtbar macht und im Austausch mit den Nutzenden schließt.
\newline	
Neben diesen inhaltlichen Anforderungen ergeben sich technische Anforderungen an Modularität und Offenheit der Systemarchitektur. Da sich sowohl die Leistungsfähigkeit von LLMs als auch die zugänglichen Modellvarianten kontinuierlich weiterentwickeln, sollte das Assistenzsystem so gestaltet sein, dass unterschiedliche Sprachmodelle flexibel eingebunden werden können. Dies betrifft sowohl gegenwärtig verfügbare lokale Modelle als auch zukünftige Modelle, deren Integration ohne grundlegende Architekturänderungen möglich sein muss. Darüber hinaus ist zu berücksichtigen, dass nicht alle Nutzenden über die technischen Voraussetzungen verfügen, ein LLM lokal auszuführen. Für solche Szenarien muss das System alternative Wege der Modellgenerierung anbieten, etwa durch die Möglichkeit, vorab generierte JSON-Strukturen einzulesen und weiterzuverarbeiten. Diese Offenheit erhöht die Nutzbarkeit des Systems unabhängig von der technischen Umgebung.

Von wissenschaftlicher Relevanz ist zudem die Fähigkeit des Systems, Prozessdaten transparent zu protokollieren. Die im Dialogverlauf entstehenden Informationen, einschließlich Rückfragen, Korrekturen und Zwischenschritte, bilden eine wertvolle empirische Grundlage für spätere Analysen, beispielsweise hinsichtlich der Qualität der Modellinterpretation oder der Interaktionsmuster. Dementsprechend sollte das System Möglichkeiten zur Erfassung und Speicherung dieser Daten bereitstellen, sei es in Form von Log-Dateien oder strukturierten JSON-Dokumenten.

Schließlich ergeben sich Anforderungen an die Gestaltung der Benutzungsoberfläche. Ein dialogbasiertes Modellierungswerkzeug muss intuitiv bedienbar sein, visuell übersichtlich bleiben und Nutzende nicht mit unnötiger Komplexität konfrontieren. Da Eingabefehler oder Fehlbedienungen den Modellierungsprozess erheblich beeinflussen können, sollte das Interface so gestaltet sein, dass Fehlinteraktionen weitgehend ausgeschlossen werden. Dies umfasst klare Rückmeldungen, konsistente Interaktionsmechanismen und eine reduzierte Komplexität der verfügbaren Aktionen.

In ihrer Gesamtheit bilden diese Anforderungen das konzeptionelle Fundament für die nachfolgenden Schritte der Datenextraktion, der formalen Modellstrukturierung und der Systemarchitektur. Sie definieren sowohl die funktionalen Erwartungen an das Assistenzsystem als auch die Rahmenbedingungen, unter denen dieses im praktischen Einsatz robust, erweiterbar und langfristig nutzbar bleibt.

\subsection*{3.3 Prinzipien der dialogbasierten Bedeutungserschließung}

Die im vorhergegangenem Kapitel formulierten Anforderungen verdeutlichen, dass ein dialogorientiertes Assistenzsystem nur dann robuste und nachvollziehbare Prozessmodelle erzeugen kann, wenn die im Gespräch geäußerten Inhalte zunächst semantisch erschlossen und inhaltlich vereinheitlicht werden. Dieser Schritt markiert den Übergang von der offenen, kontextabhängigen Struktur natürlicher Sprache hin zu einer abstrahierten Bedeutungsform, die einer strukturierten Modellbildung zugänglich ist. Die dialogische Interaktion bildet dabei nicht lediglich ein Bedienprinzip, sondern eine methodische Grundlage, um Unschärfen, implizite Bedeutungen und fehlende Informationen systematisch zu erkennen und zu beheben.

Die natürliche Sprache unterliegt den Regeln menschlicher Kommunikation, die sich durch Variabilität, Fragmentierung und kontextabhängige Bedeutungsbildung auszeichnet. Prozessbeschreibungen entstehen selten als vollständig formulierte, logisch strukturierte Darstellungen, sondern entwickeln sich schrittweise im Gespräch. Die dialoggestützte Interpretation besteht daher nicht darin, einzelne Aussagen isoliert zu interpretieren, sondern darin, aus dem fortlaufenden Austausch eine logisch aufgebaute Prozesslogik zu rekonstruieren. Moderne Sprachmodelle können diesen Prozess unterstützen, indem sie semantische Muster, Rollenbezüge und Bedingungen identifizieren, benötigen jedoch zusätzliche Mechanismen zur Sicherstellung von Konsistenz, Vollständigkeit und formaler Interpretierbarkeit.


\subsubsection*{3.3.1 Charakteristika natürlichsprachlicher Prozessbeschreibungen}

Natürlichsprachliche Prozessbeschreibungen weisen mehrere strukturelle Eigenschaften auf, die ihre direkte Überführung in formale Modelle erschweren. Erstens sind sie häufig fragmentiert. Informationen zu Aktivitäten, Entscheidungen oder Rollen werden über mehrere Gesprächsbeiträge verteilt, teilweise in unterschiedlicher Reihenfolge oder in unterschiedlichen Detailgraden. 
Zweitens enthalten sprachliche Äußerungen oftmals implizite Bedeutungen, die sich erst aus dem Kontext oder aus dem Hintergrundwissen der Sprechenden ergeben. Formulierungen wie „danach prüfen wir das“ oder „falls das nicht geht“ enthalten nur dann prozessrelevante Bedeutung, wenn der vorherige Gesprächsverlauf berücksichtigt wird.
Drittens weisen Prozessbeschreibungen eine ausgeprägte Kontextabhängigkeit auf. Die Interpretation einer Aussage hängt stark von vorherigen Beiträgen, unausgesprochenen Annahmen oder organisationalem Wissen ab. 
Viertens findet häufig eine semantische Verdichtung statt:
Mehrere prozessrelevante Aspekte, etwa Bedingungen, Rollen, Auslöser oder Folgeaktivitäten, werden in einem einzelnen Satz kodiert. Der Satz „Wenn die Unterlagen fehlen, informieren wir den Kunden“ enthält beispielsweise eine Bedingung, eine Entscheidung, eine Rolle und eine Folgehandlung.

Diese Eigenschaften machen deutlich, dass ein Assistenzsystem nicht auf eine lineare Extraktion einzelner Informationen zurückgreifen kann. Vielmehr muss es Mechanismen bereitstellen, die dialogische Kohärenz herstellen, Kontext berücksichtigen und Bedeutungen aus diesem heraus argumentativ entwickeln. Erst diese semantische Vorstrukturierung ermöglicht eine spätere formale Modellbildung.


\subsubsection*{3.3.2 Dialogische Bedeutungserschließung durch Sprachmodelle}

Große Sprachmodelle eignen sich gut, um natürlichsprachliche Prozessbeschreibungen zu verstehen, weil sie Bedeutungen im Kontext erfassen können. Sie erkennen Zusammenhänge über mehrere Gesprächsbeiträge hinweg, füllen fehlende Informationen durch Schlussfolgerungen auf und identifizieren Handlungsmuster, die für die Prozessmodellierung wichtig sind. Gleichzeitig arbeiten LLMs jedoch nicht vollständig vorhersehbar und können zu Unsicherheiten, Mehrdeutigkeiten oder Halluzinationen führen. Ein rein automatisches Übersetzen von natürlicher Sprache in formale Prozesslogiken wäre deshalb nicht zuverlässig.

Daher spielt die dialogische Interaktion eine zentrale Rolle. Das Assistenzsystem muss erkennen, wenn etwas im Sprachinput unklar ist, und gezielt Rückfragen stellen. Diese dienen nicht nur dazu, fehlende Informationen zu ergänzen, sondern sorgen vor allem für eine präzise Bedeutungsbestimmung. Die wiederkehrende Abstimmung zwischen Nutzenden und System folgt dem in Kapitel 2.3 beschriebenen dialogischen Kommunikationsmodell und bildet die Grundlage für eine gemeinsame Modellierung, in der Prozesswissen Schritt für Schritt konkretisiert und verfeinert wird.

Das Zusammenspiel aus LLM-basierter Bedeutungsinterpretation und dialogischer Klärung ermöglicht es, aus offenen, variablen und teilweise unstrukturierten sprachlichen Äußerungen konsistente semantische Einheiten zu formen. Diese bilden die Grundlage für die nachfolgende Strukturierung in ein semiformalen Zwischenmodell.


\subsection*{3.4 Semiformale Repräsentation und Modellbildung}

Die im Dialogverlauf gewonnenen Strukturen müssen in eine Form überführt werden, die sowohl für technische Komponenten interpretierbar als auch für die spätere BPMN-Generierung hinreichend formalisiert ist. Eine solche Zwischenrepräsentation fungiert als Bindeglied zwischen den aus der Sprache erschlossenen semantischen Einheiten und der formalen Modelllogik der BPMN. Sie bildet zentrale prozessuale Entitäten – etwa Aktivitäten, Ereignisse, Rollen und Kontrollflussbeziehungen – eindeutig und strukturiert ab, ohne dabei die vollständige Komplexität der grafischen Notation vorwegzunehmen.

Aus theoretischer Perspektive erfüllt ein solches Zwischenmodell eine doppelte Funktion. Erstens ermöglicht es die Abstraktion natürlicher Sprache in eine formalere Ebene, auf der prozessrelevante Informationen explizit dargestellt werden. Zweitens bildet es die Grundlage für die Modellbildung, bei der aus diesen strukturierten Informationen ein formal gültiges BPMN-Modell abgeleitet wird. Entsprechend muss ein Zwischenmodell Anforderungen wie Vollständigkeit, Eindeutigkeit, Erweiterbarkeit, Fehlertoleranz und eine klar interpretierbare Struktur erfüllen.

In dieser Arbeit wird die Zwischenrepräsentation als semiformale Datenstruktur konzipiert, deren technische Umsetzung im Prototyp auf JSON basiert. Diese Entscheidung folgt nicht technischen Zwängen, sondern ergibt sich aus der Notwendigkeit einer hierarchischen, flexibel erweiterbaren und leicht validierbaren Struktur. Die JSON-basierte Repräsentation ist somit Ausdruck eines konzeptionellen Gestaltungsprinzips. Sie erlaubt es, Rollenbezüge, Ablaufstrukturen und logische Beziehungen präzise zu definieren und zugleich den Dialogverlauf sowie Import- und Exportmechanismen zu unterstützen.

Auf dieser Grundlage erfolgt im nächsten Schritt die Modellbildung, in der die semantischen Einheiten in BPMN-Elemente überführt und in einer formal zulässigen Struktur angeordnet werden. Dieser Prozess umfasst insbesondere die Ableitung von Aufgaben, Ereignissen und Gateways, die Spezifikation zulässiger Kontrollflüsse sowie die Sicherstellung syntaktischer und semantischer Konsistenz im Sinne der in Kapitel 2.1 beschriebenen BPMN-Regelwerke. Die Modellbildung schließt die konzeptionelle Phase ab und bereitet die Grundlage für die anschließende technische Umsetzung im Prototyp.

\subsection*{3.5 Konzeptionelle Systemarchitektur}

Die zuvor beschriebenen Anforderungen an die Interpretation, Strukturierung und Validierung natürlicher Sprache lassen sich nur durch ein mehrschichtiges Systemdesign erfüllen, das unterschiedliche funktionale Perspektiven trennt und zugleich einen durchgängigen Informationsfluss sicherstellt. Die Architektur folgt daher einem modularen Ansatz, in dem dialogische, semantische, strukturelle und technische Komponenten klar abgegrenzt sind.

Im Zentrum der Architektur steht eine Interaktionsebene, die den dialogischen Prozess steuert. Sie verwaltet den Gesprächskontext, stellt dem Sprachmodell relevante historische Informationen bereit und vermittelt zwischen Nutzenden und semantischer Analyse. Dadurch wird gewährleistet, dass modellierungsrelevante Aussagen nicht isoliert, sondern im Verlauf des gesamten Dialogs interpretiert werden. Gleichzeitig übernimmt diese Ebene eine Kontrollfunktion, indem sie Rückfragen generiert, sobald Unsicherheiten oder strukturelle Lücken erkennbar sind.

Darunter liegt eine semantische Verarbeitungsebene, die die natürlichsprachlichen Eingaben interpretiert und in domänenspezifische Bedeutungsstrukturen überführt. Sie nutzt die Fähigkeiten großer Sprachmodelle zur kontextsensitiven Bedeutungserschließung, strukturiert deren Output jedoch gezielt, indem sie ihn auf prozessrelevante Entitäten wie Aktivitäten, Ereignisse, Rollen und Kontrollflusslogiken reduziert. Damit entsteht die konzeptionelle Grundlage für den Übergang zu formalen Modellstrukturen.

Die Repräsentationsebene überführt diese Bedeutungsstrukturen in eine wohldefinierte Zwischenform, die als Bindeglied zwischen Sprachverarbeitung und formaler Modellbildung dient. Die in dieser Arbeit gewählte JSON-Struktur fungiert als semiformales Modell, das sowohl für Sprachmodelle gut erzeugbar als auch für nachgelagerte Komponenten eindeutig interpretierbar ist. Sie stellt das strukturelle Gerüst bereit, auf dem später ein vollständiges BPMN-Modell aufgebaut werden kann.

Eine darauf aufbauende Validierungsebene prüft die Zwischenrepräsentation auf interne Konsistenz, Vollständigkeit und semantische Plausibilität. Sie erkennt beispielsweise fehlende Verbindungen, widersprüchliche Rollen oder unzulässige Kontrollflussstrukturen. Werden solche Probleme festgestellt, werden sie an die Interaktionsebene zurückgegeben, wodurch ein iterativer, dialogisch eingebetteter Modellierungsprozess entsteht.

Die Modellbildungsebene definiert anschließend die abstrakte Struktur des späteren BPMN-Modells. Sie erzeugt die logischen Beziehungen zwischen den Modell­elementen und bereitet deren Anordnung vor, ohne jedoch konkrete grafische Layoutentscheidungen zu treffen.

Obwohl die prototypische Umsetzung im weiteren Verlauf dieser Arbeit auf spezifische Technologien zurückgreift, etwa ein in Godot eingebettetes UI-System, ein lokales LLM über Ollama oder modularisierte Renderingkomponenten –, beschreibt die hier dargestellte Architektur eine davon abstrahierte konzeptionelle Struktur. Sie zeigt, welche funktionalen Rollen das spätere System erfüllen muss und wie sich dialogisch erhobene Prozessbeschreibungen systematisch in formal definierte Modellstrukturen überführen lassen.

\section*{4. Prototypische Implementierung}

Auf Grundlage des in Kapitel~3 entwickelten Konzepts wurde ein vollständig funktionsfähiger Prototyp umgesetzt, der den dialogbasierten Übergang von natürlichsprachlichen Prozessbeschreibungen zu formalisierten BPMN-Modellen realisiert. Während das vorhergehende Kapitel die theoretische und methodische Fundierung des Ansatzes gelegt hat, beschreibt das nun folgende Kapitel die konkrete technische Ausgestaltung des Systems. Die Implementierung folgt dabei einem modularen Aufbau, der Sprachverarbeitung, Modellgenerierung, Layoutberechnung und visuelle Darstellung klar voneinander trennt und zugleich in ein übergreifendes Interaktionsmodell einbettet. Der Prototyp stellt damit nicht nur einen Machbarkeitsnachweis dar, sondern veranschaulicht, wie sich die theoretischen Überlegungen praktisch in einer lauffähigen Systemarchitektur umsetzen lassen.

\subsection*{4.1 Technologiestack und Systemumgebung}

Die technische Realisierung des Prototyps erfolgte in einer vollständig eigenständigen Systemumgebung, die gezielt darauf ausgelegt ist, die Interaktion zwischen Benutzereingaben, Large Language Model und BPMN-Visualisierung möglichst eng zu verzahnen. Die Entscheidung für die verwendeten Technologien orientierte sich sowohl an funktionalen Anforderungen aus Kapitel~3 als auch an pragmatischen Überlegungen hinsichtlich Performanz, Erweiterbarkeit und lokaler Ausführbarkeit.

\subsubsection*{4.1.1 Godot Engine und GDScript}

Die Implementierung erfolgte in der Godot Engine in der Version 4.5.1, die aufgrund ihrer Flexibilität, ihrer leistungsfähigen 2D-Rendering-Architektur und ihrer szenenbasierten Struktur besonders geeignet ist, komplexe visuelle Modelle wie BPMN-Diagramme interaktiv darzustellen. Die Engine bietet eine klare Trennung zwischen Logik, Darstellung und Datenstrukturen, was für die Abbildung der Modellierungs- und Layoutschritte von zentraler Bedeutung ist. Die Skripterstellung erfolgte in GDScript, einer Python-ähnlichen, leichtgewichtigen Sprache, die eine hohe Lesbarkeit sowie eine enge Integration mit den internen Engine-Mechanismen ermöglicht. 

Godot erwies sich zudem als vorteilhaft für die Umsetzung dynamischer Interfaces, wie sie für dialogbasierte Assistenzsysteme erforderlich sind. Das LlmChatWindow bildet in diesem Zusammenhang die zentrale Interaktionsoberfläche, 
in der Benutzende ihre Prozessbeschreibungen eingeben und Rückmeldungen des Systems abrufen können. Die Wahl einer Game-Engine mag auf den ersten Blick untypisch für eine klassische BPMN-Anwendung erscheinen; im Kontext dieser Arbeit ermöglicht sie jedoch eine besonders intuitive grafische Darstellung, flüssige Animationen und die leichte Erweiterbarkeit um interaktive Elemente.

\subsubsection*{4.1.2 LLM-Ausführung mit Ollama}

Für die Sprachverarbeitung wird ein lokal ausgeführtes LLM über die Plattform Ollama genutzt. Die Entscheidung für eine lokale Modellinstanz basiert auf zwei zentrale Erwägungen: Erstens ermöglicht sie die Verarbeitung sensibler Prozessinformationen ohne externe Cloud-Dienste und trägt damit den Anforderungen an Datenschutz und Nachvollziehbarkeit Rechnung. Zweitens erlaubt die lokale Ausführung eine präzise Kontrolle über das verwendete Modell, seine Konfigurationen und die vollständige Reproduzierbarkeit der Ergebnisse.

Die Kommunikation zwischen Anwendung und LLM erfolgt über den eigens entwickelten OllamaClient, der über eine HTTP-basierte Schnittstelle in den ChatController eingebunden ist. Dieser steuert die Generierung, verwaltet den Dialogkontext und sorgt dafür, dass der systemseitige Prompt aus Kapitel~3 strikt eingehalten wird. Durch die Nutzung eines generativen Modells ist es möglich, komplexe Prozessbeschreibungen kontextsensitiv zu interpretieren und sowohl strukturierte JSON-Modelle als auch Rückfragen bei Unklarheiten zu erzeugen.

\subsubsection*{4.1.3 Modellarchitektur des verwendeten Sprachmodells}

Das verwendete Sprachmodell basiert auf der LLaMA-3-Architektur, die sich durch eine hohe Effizienz und starke Fähigkeiten im Bereich der strukturierten Textgenerierung auszeichnet. Die Parametergröße ist so gewählt, dass einerseits ausreichende Komplexität für die Interpretation von Prozesslogik verfügbar ist, andererseits eine performante Ausführung auf handelsüblicher Hardware gewährleistet bleibt. 

Im Rahmen des Prototyps kommt das Modell primär für drei Aufgaben zum Einsatz:
(1) die semantische Segmentierung und Interpretation natürlichsprachlicher Eingaben,  
(2) die Transformation der Extraktionsergebnisse in eine strukturierte BPMN-JSON-Repräsentation  
und (3) die dialogische Klärung von Ambiguitäten.  

Der Masterprompt, der den gesamten Modellierungsprozess steuert, definiert strikte syntaktische Regeln, eine feste Elementkategorie nach BPMN-Standard sowie Validierungsmechanismen, die vom Modell einzuhalten sind. Damit wird die 
potentielle Varianz generativer Modelle gezielt eingeschränkt, um eine hohe Strukturqualität der erzeugten Modelle sicherzustellen.

\subsubsection*{4.1.4 Frontend- und Renderkomponenten}

Die visuelle Darstellung der BPMN-Modelle erfolgt vollständig innerhalb der Godot-Engine und basiert auf einer Sammlung spezialisierter 2D-Szenen, die den grafischen Elementen der BPMN entsprechen. Ereignisse, Aktivitäten und Gateways werden jeweils durch eigens entwickelte Node-Typen repräsentiert, deren visuelle Struktur den Notationselementen der BPMN-2.0-Spezifikation folgt. 

Das Frontend umfasst darüber hinaus ein interaktives Dialogfenster, in dem sowohl die Eingaben der Nutzenden als auch die Antworten des Systems dargestellt werden. Dieses Interface verbindet die textbasierte Interaktion direkt mit dem Modellierungsprozess und ermöglicht es, den Prozessablauf Schritt für Schritt nachzuvollziehen. Die Integration der visuellen Modellkomponenten in die dialogbasierte Umgebung unterstützt somit ein essentielles Ziel aus Kapitel~2: die direkte Kopplung von natürlicher Interaktion und formalisierter Prozessdarstellung.

\subsection*{4.2 Implementierung des LLM-basierten BPMN-Generators}

Die Generierung strukturierter BPMN-Modelle aus natürlichsprachlichen Beschreibungen stellt den Kern des entwickelten Assistenzsystems dar. Aufbauend auf den in Kapitel~3 entwickelten konzeptionellen Grundlagen wurde ein mehrstufiger Generierungsprozess implementiert, der das Large Language Model nicht lediglich als Textgenerator verwendet, sondern als zentrale semantische Instanz, die in der Lage ist, Prozesslogiken zu interpretieren, strukturelle Zusammenhänge abzuleiten und das resultierende Modell in einer formalisierten JSON-Struktur darzustellen. Die Implementierung folgt dabei einem kontrollierten und stark regelbasierten Ansatz, der die generative Offenheit des Modells begrenzt und gleichzeitig seine interpretativen Fähigkeiten nutzbar macht.

\subsubsection*{4.2.1 Prompt-Design und Dialogsteuerung}

Der gesamte Generierungsprozess wird durch einen ausführlich ausgearbeiteten Masterprompt gesteuert, der dem LLM als kontextuelle Leitlinie dient. Die Ausarbeitung dieses Prompts orientiert sich an den theoretischen Erkenntnissen zu Prompt Engineering aus Kapitel~2.2.2 und bildet eine der zentralen technischen Stellgrößen des Systems.

Der Prompt erfüllt mehrere Funktionen zugleich:  
Er definiert erstens die Rolle des Modells als strikt regelbefolgenden BPMN-Assistenten, der aus beliebig formulierten Beschreibungen ausschließlich strukturierte JSON-Ausgaben erzeugen darf. Zweitens enthält er präzise Vorgaben zu den erlaubten Elementtypen, der zulässigen Struktur der Sequenzflüsse und der korrekten ID-Vergabe. Drittens beschreibt er Validierungsschritte, die das Modell vor der Ausgabe intern durchzuführen hat, sowie die Anforderung, bei Unklarheiten gezielte Rückfragen zu formulieren. 

Die Dialogsteuerung erfolgt über den ChatController, der die Nachrichtenhistorie verwaltet und dafür sorgt, dass das Modell stets mit einem konsistenten Kontext arbeitet. Dabei wird bewusst nur ein reduzierter Ausschnitt des bisherigen Dialogs übermittelt, um Halluzinationen und Kontextverwässerungen zu minimieren. Der Controller fungiert gleichzeitig als Kommunikationsorgan zwischen Nutzer, Systemlogik und LLM, sodass der Transformationsprozess als fortlaufender Dialog gestaltet wird, der semantische Korrektheit und Interaktivität vereint.

\subsubsection*{4.2.2 JSON-Validierung und Parsing (BpmnJsonLoader)}

Das vom LLM erzeugte JSON bildet die strukturierte Zwischenrepräsentation des Prozessmodells. Da generative Modelle trotz starker Regelvorgaben fehleranfällig bleiben können, wurde ein Parser implementiert, der die erzeugten Daten systematisch prüft und in die internen BPMN-Datenstrukturen des Prototyps überführt. Die Validierung umfasst sowohl syntaktische als auch semantische Aspekte und stellt sicher, dass die Ausgabe des LLMs den Anforderungen 
des BPMN-Standards entspricht.

Der BpmnJsonLoader übernimmt dabei mehrere Aufgaben. Zunächst wird die JSON-Struktur auf Vollständigkeit geprüft, insbesondere hinsichtlich der Element-IDs, der Sequenzflüsse und der Beschriftungen. Anschließend erfolgt die  Übertragung der Daten in Godot-spezifische Instanzen, die den späteren Visualisierungsprozess ermöglichen. Dabei entstehen konkrete Objekte wie TaskNode2D, StartEvent2D oder ExclusiveGateway2D, die später von den Layout- und  Renderkomponenten weiterverarbeitet werden.

Ein besonderer Schwerpunkt liegt auf der Prüfung der Flussbeziehungen. Da Prozessmodelle häufig durch implizite Logik beschrieben werden, kann es zu Abweichungen in den Referenzen oder der Strukturierung kommen. Der Parser gleicht daher alle referenzierten IDs ab, stellt die Existenz der Zielknoten sicher und markiert Inkonsistenzen, die in einem separaten Korrekturschritt behoben werden können. Die JSON-Verarbeitung bildet damit den entscheidenden Übergang zwischen der sprachlichen Interpretation des LLMs und der graphischen Darstellung innerhalb der Engine und gewährleistet, dass nur formal gültige Strukturen an den Modellierungsprozess übergeben werden.

\subsubsection*{4.2.3 Fehlerbehandlung und Autokorrekturmechanismen}

Trotz klarer Vorgaben neigen generative Modelle dazu, gelegentlich syntaktische oder strukturelle Fehler zu erzeugen, etwa durch fehlende Referenzen, unvollständige Elementangaben oder unzulässige Sequenzflüsse. Um diesen Herausforderungen zu begegnen, wurde ein Mechanismus implementiert, der sowohl auf der Ebene des Prompts als auch auf der Ebene der Systemlogik greift.

Auf Ebene des LLMs erzwingt der Masterprompt eine interne Vorvalidierung, bei der das Modell überprüft, ob alle Elemente konsistent strukturiert sind. Sollte es Unstimmigkeiten erkennen, muss es vor der finalen Ausgabe einen Fehlerbericht erzeugen und anschließend eine korrigierte Struktur liefern. Dieser Ansatz nutzt bewusst die Fähigkeit moderner LLMs, eigene Ausgaben zu reflektieren und zu reparieren.

Auf Systemebene unterstützt der Parser diesen Prozess, indem er Fehlerkennungen protokolliert und gegebenenfalls Rückmeldungen an das LLM initiieren kann. In der praktischen Umsetzung zeigte sich, dass diese Kombination aus  modellinterner Selbstkorrektur und externer Validierung zu einer deutlich höheren Strukturqualität führt. Der Generator entwickelt sich damit zu einem hybriden Regel- und Sprachmodell, das den Balancepunkt zwischen generativer Flexibilität und struktureller Strenge wahrt.

Die Fehlerbehandlung trägt schließlich dazu bei, dass Nutzerinnen und Nutzer auch bei komplexen oder unvollständigen Eingaben nachvollziehbare und konsistente Modelle erhalten. Gleichzeitig schafft sie die notwendige Transparenz, die für vertrauenswürdige KI-gestützte Assistenzsysteme im Sinne der in Kapitel~2 erörterten wertegeleiteten Interaktionen essenziell ist.

\subsection*{4.3 BPMN-Layout und visuelle Darstellung}

Nachdem die durch das LLM erzeugte JSON-Struktur in interne BPMN-Objekte überführt wurde, besteht der nächste wesentliche Schritt in der räumlichen Anordnung und anschließenden graphischen Darstellung der Modellelemente. Die visuelle Aufbereitung ist nicht lediglich ein kosmetischer Bestandteil des Systems, sondern ein integraler Teil der Modellinterpretation, da erst durch die räumliche Strukturierung die inhaltlichen Zusammenhänge, Kontrollflüsse 
und Prozesslogiken wahrnehmbar werden. Die Gestaltung eines angemessenen Layouts stellt daher eine eigenständige konzeptionelle und technische Herausforderung dar, insbesondere vor dem Hintergrund, dass die Modelle aus variantenreichen und unstrukturierten sprachlichen Beschreibungen generiert werden.

Um dieser Aufgabe gerecht zu werden, wurde eine modulare Layoutpipeline entwickelt, die mehrere spezialisierte Komponenten umfasst. Diese trennt den Prozess der räumlichen Positionierung strikt von der semantischen Modellgenerierung und ermöglicht es, die algorithmische Logik vollständig unabhängig von der Sprachverarbeitung zu halten. Dadurch entsteht ein System, das die in Kapitel~3 beschriebenen Anforderungen an Transparenz, Wiederholbarkeit und Strukturtreue erfüllt.

\subsubsection*{4.3.1 Architektur der Layout-Pipeline (LayoutHandler)}

Im Zentrum der Layoutverarbeitung steht der eigens entwickelte LayoutHandler, der die gesamte Pipeline koordiniert und als übergeordnete Instanz fungiert, welche die repräsentierten BPMN-Elemente einliest, strukturiert und an die Layoutengine weiterreicht. Der Handler interpretiert die Sequenzflüsse als gerichteten Graphen, dessen Struktur die Grundlage für die spätere räumliche Anordnung bildet. Dieser Graph wird anschließend in eine hierarchische Sequenz überführt, die Start- und Endereignisse, Zwischenknoten, Verzweigungen und Zusammenführungen eindeutig ordnet. 

In dieser Phase wird nicht versucht, die exakte visuelle Position einzelner Elemente festzulegen; vielmehr erzeugt der LayoutHandler eine abstrakte Prozessstruktur, die die logische Reihenfolge und Gruppierung der BPMN-Elemente widerspiegelt. Dieser Zwischenschritt ist notwendig, um die Heterogenität natürlicher Sprache zu harmonisieren und Fehler, die durch inkonsistente oder unvollständige Eingaben entstehen können, abzufangen. Erst aus dieser bereinigten Struktur kann eine stabile visuelle Darstellung abgeleitet werden.

\subsubsection*{4.3.2 Layoutalgorithmen (LayoutEngine)}

Die eigentliche Positionierungslogik wird durch die LayoutEngine umgesetzt, die auf graphbasierten und rasterorientierten Verfahren beruht. Ziel dieser Komponente ist es, eine klare, übersichtliche und semantisch interpretierbare Anordnung zu erzeugen, die sowohl den Konventionen der BPMN-Notation als auch den visuellen Erwartungen an Geschäftsprozessdiagramme entspricht.

Die Engine analysiert die vom LayoutHandler bereitgestellte abstrakte Struktur und ordnet die Elemente zunächst entlang eines sequentiellen Flusspfads an, wobei Verzweigungen durch horizontale Abspaltungen und parallele Abläufe durch vertikale Ausdifferenzierungen dargestellt werden. Diese horizontale–vertikale Gliederung folgt den in der Modellierungspraxis üblichen Konventionen und erleichtert die spätere Interpretation des Diagramms.

Ein besonderes Augenmerk gilt der Behandlung komplexer Gateway-Strukturen. Entscheidungen und parallele Pfade werden so positioniert, dass Ein- und Ausgänge klar sichtbar bleiben und keine Überlagerungen entstehen. Die Engine 
berücksichtigt dabei geometrische Randbedingungen wie Mindestabstände, Beschriftungsflächen und Assoziationswinkel. Dadurch wird verhindert, dass der Prozessfluss visuell verwischt oder dass semantische Beziehungen  missverständlich dargestellt werden.

\subsubsection*{4.3.3 Routing der Sequenzflüsse (FlowConnectionHandler, AutoConnector)}

Sobald die Positionen der Modellelemente festgelegt sind, werden die Verbindungen zwischen ihnen berechnet. Diese Aufgabe übernehmen der FlowConnectionHandler und der AutoConnector, zwei Komponenten, die gemeinsam die grafische Umsetzung der BPMN-Sequenzflüsse realisieren.

Der FlowConnectionHandler analysiert zunächst die logischen Beziehungen zwischen den Knoten, wie sie aus der JSON-Struktur hervorgehen, und erzeugt eine interne Repräsentation der Flusskanten. Dabei werden auch semantische  Besonderheiten berücksichtigt, etwa Gateway-Ausgänge mit spezifischen Bedingungslabels oder parallele Flüsse, die synchron beginnen oder enden.

Der AutoConnector übernimmt anschließend die geometrische Umsetzung dieser Flusskanten. Er berechnet orthogonale Routen, die sich an gängigen Diagrammkonventionen orientieren und sowohl die Übersichtlichkeit als auch die  Lesbarkeit steigern. Besondere Herausforderungen ergeben sich bei kreuzenden oder stark verschachtelten Kanten, die durch adaptive Routing-Algorithmen vermieden oder visuell entschärft werden. Der AutoConnector berücksichtigt  zudem die Ankerpunkte der BPMN-Elemente, sodass ein natürlicher und ästhetischer Fluss entsteht, der den Wahrnehmungsgewohnheiten in der Prozessmodellierung entspricht.

\subsubsection*{4.3.4 Renderer für BPMN-Elemente}

Die abschließende Phase der visuellen Darstellung wird durch den Renderer umgesetzt, der die zuvor angeordneten und verbundenen Elemente in Godot als vollwertiges BPMN-Diagramm ausgibt. Jedes Element wird durch eine spezialisierte Szene repräsentiert, die die grafische Notation exakt nach den Spezifikationen der BPMN-2.0-Standards abbildet. Dies umfasst sowohl die Grundformen – etwa Kreise für Ereignisse, Rauten für Gateways oder 
abgerundete Rechtecke für Aufgaben – als auch interne Beschriftungen und Symbole.

Die Wahl einer szenenbasierten Architektur ermöglicht es, auch komplexe Modelle modular und effizient darzustellen. Änderungen im Layout oder in den Verbindungen werden automatisch propagiert, sodass das Diagramm stets konsistent und vollständig bleibt. In Kombination mit den adaptiven Layoutkomponenten entsteht so ein dynamisches Visualisierungssystem, das die durch das LLM erzeugten Prozessstrukturen nicht nur korrekt abbildet, sondern auch für Nutzende intuitiv nachvollziehbar macht.

Mit der Fertigstellung der Layout- und Renderstufe ist der Generierungsprozess abschließend vollständig visualisiert. Die resultierende Darstellung bildet die Grundlage für die nachfolgende Integration in das dialogbasierte 
Assistenzsystem, das in Abschnitt~4.4 beschrieben wird.

\subsection*{4.4 Integration in ein dialogbasiertes Assistenzsystem}

Die zuvor beschriebenen Komponenten zur Sprachverarbeitung, Modellgenerierung und Layoutberechnung entfalten ihr volles Potenzial erst in Kombination mit einer interaktiven Benutzerschnittstelle, die den dialogischen Charakter des Systems erlebbar macht. Die Integration dieser Module in ein konsistentes Assistenzsystem stellt somit einen zentralen Schritt der prototypischen Umsetzung dar. Ziel ist es, eine Umgebung zu schaffen, in der Nutzende ihre Prozessbeschreibungen in natürlicher Sprache formulieren können und die resultierenden BPMN-Modelle in Echtzeit nachvollziehbar generiert werden. Die Interaktion mit dem System bildet damit den verbindenden Rahmen zwischen Spracheingabe, Modelltransformation und visualisiertem Ergebnis.

\subsubsection*{4.4.1 Chatinterface (LlmChatWindow)}

Das Chatinterface stellt die primäre Interaktionsfläche des Prototyps dar. Im LlmChatWindow Items angezeigt, wodurch sich ein dialogischer Fluss ergibt, der den Strukturierungsprozess unmittelbar transparent macht. Das Interface ist so gestaltet, dass es einer linearen Chat-Timeline ähnelt, in der Konversationen natürlich nachvollziehbar dargestellt werden. Dadurch wird die kognitive Brücke zwischen freier Spracheingabe und formalem Modell erleichtert. 

Ein wesentliches Designprinzip besteht darin, die Komplexität technischer Abläufe vor den Nutzenden zu abstrahieren, ohne den Prozess intransparenter zu machen. So werden beispielsweise Fehlerhinweise oder Rückfragen des Systems direkt im Kontext des Gesprächs angezeigt, was die Orientierung fördert und dem wertegeleiteten Interaktionsprinzip aus Kapitel~2.3 Rechnung trägt.

Das Chatinterface fungiert zudem als Bindeglied zwischen Benutzerinput und ChatController. Es leitet Eingaben weiter, löst Generierungsprozesse aus und aktualisiert die Darstellung der Systemantworten dynamisch. Damit bildet es 
den zentralen operativen Zugangspunkt zum Assistenzsystem.

\subsubsection*{4.4.2 Interaktionslogik und Zustandsverwaltung}

Die dialogische Struktur des Systems setzt eine robuste Interaktionslogik voraus, die sicherstellt, dass Nutzereingaben korrekt verarbeitet, Sequenzen nachvollzogen und Modellkontexte konsistent gehalten werden. Diese Aufgaben werden im Kern durch den ChatController übernommen, der als orchestrierende Logikschicht fungiert. Er verwaltet die Nachrichtenhistorie, kontrolliert die Kommunikation mit dem LLM und stellt sicher, dass der systemseitige Prompt stets vollständig berücksichtigt wird.

Ein besonderes Augenmerk liegt auf der Kontextverwaltung. Sprachliche Prozessbeschreibungen erfolgen selten vollständig in einer einzelnen Äußerung; vielmehr ergeben sie sich aus einer Reihe aufeinander aufbauender  Dialogbeiträge. Der ChatController implementiert daher Mechanismen zur kontrollierten Kontextbegrenzung, sodass das Modell nur diejenigen Teile der Historie erhält, die für die Interpretation der aktuellen Eingabe relevant sind. Diese Maßnahme reduziert nicht nur die Rechenlast, sondern minimiert auch das Risiko von fehlerhaften oder übermäßig kreativen Modellgenerierungen durch sogenannte LLM-Halluzinationen (vgl.~Kap.~2.2.3).

Die Zustandsverwaltung umfasst darüber hinaus die Koordination der Übergänge zwischen Generierungsphasen. Nach Eingang der LLM-Ausgabe löst der Controller den Parsingprozess aus, initiiert die Layoutpipeline und aktualisiert schließlich die Visualisierung. Dadurch entsteht ein in sich geschlossener Kreislauf, der die Transformation von Sprache zu Diagramm konsistent steuert.

\subsubsection*{4.4.3 Nutzerführung und Rückfragenmechanismen}

Ein zentrales Merkmal der prototypischen Umsetzung besteht darin, dass Unklarheiten in den Nutzereingaben nicht zu Fehlgenerierungen führen, sondern explizit durch Rückfragen adressiert werden. Diese Fähigkeit basiert auf Eigenschaften moderner LLMs, die es ermöglichen, semantische Lücken oder Ambiguitäten im Gesprächskontext zu identifizieren und gezielt nachzufragen. Dies entspricht nicht nur interaktionspsychologischen Prinzipien, sondern auch den Anforderungen an vertrauenswürdige Assistenzsysteme, wie sie in Kapitel~2.3 herausgearbeitet wurden.

In der praktischen Umsetzung wird dieses Verhalten durch die Promptlogik gesteuert. Das Modell ist instruiert, nur dann ein BPMN-Modell auszugeben, wenn alle notwendigen Strukturinformationen vorliegen. Sollte dies nicht der Fall sein, formuliert es eine präzise Rückfrage, die auf die fehlende Information verweist. Dieser Mechanismus reduziert nicht nur Fehler in der Modelltransformation, sondern unterstützt auch den iterativen Charakter des Modellierungsprozesses, indem Nutzende schrittweise zur Präzisierung ihrer Beschreibungen angeleitet werden.

Die Nutzerführung erstreckt sich zudem auf visuelle Rückmeldungen. Sobald ein Modell erfolgreich generiert wurde, erscheint die graphische Darstellung direkt neben der Konversation, wodurch eine unmittelbare Rückkopplung zwischen 
sprachlicher Eingabe und visueller Modellrepräsentation entsteht. Dieser zweikanalige Interaktionsansatz erleichtert es den Nutzenden, ihre Prozesslogik zu verstehen, zu validieren und gegebenenfalls im nächsten Dialogschritt zu korrigieren.

Insgesamt schafft das dialogbasierte Assistenzsystem eine kohärente und intuitive Benutzererfahrung, die natürliche Sprache als Ausgangspunkt nutzt und deren Transformation in eine formale BPMN-Notation transparent und nachvollziehbar gestaltet. Damit knüpft die Umsetzung nicht nur an die konzeptionellen Grundlagen aus Kapitel~3 an, sondern demonstriert auch praktisch, wie LLM-basierte Systeme in modellierenden Tätigkeiten unterstützen können.

\subsection*{4.5 Beispielhafte Konversationen und Generierungsergebnisse}

Um die Funktionsweise des entwickelten Prototyps anschaulich darzustellen und die Integration der zuvor beschriebenen Komponenten im praktischen Einsatz zu demonstrieren, werden im Folgenden exemplarische Dialoge zwischen Nutzenden und dem Assistenzsystem vorgestellt. Die Beispiele illustrieren unterschiedliche Prozessstrukturen und zeigen, wie natürliche Sprache mittels des beschriebenen LLM-basierten Generators in formal gültige BPMN-Modelle transformiert wird. In diesem Zusammenhang dienen die Konversationen nicht nur als Nachweis der technischen Funktionsfähigkeit, sondern auch als qualitative Indikatoren dafür, inwiefern der Ansatz den Anforderungen aus 
Kapitel~3 gerecht wird.

Die Beispiele wurden aus realen Interaktionen mit dem entwickelten System entnommen und sind repräsentativ für typische Nutzerszenarien. Sie spiegeln insbesondere die dialogische Natur des Ansatzes wider, da das LLM bei Bedarf Rückfragen stellt und somit aktiv zur Klärung von Unschärfen im Prozessverständnis beiträgt. Auf diese Weise entsteht ein iterativer Modellierungsprozess, der Elemente traditioneller Interaktionssysteme mit den semantischen Fähigkeiten moderner Sprachmodelle verbindet.

\subsubsection*{4.5.1 Beispiel 1: Linearer Prozess}

Linear strukturierte Prozesse bieten einen geeigneten Ausgangspunkt, um die Grundfunktion des Systems darzustellen. In einem exemplarischen Dialog beschreibt eine Nutzerin den einfachen Ablauf eines Anfrageprozesses:

\begin{quote}
	\textit{Nutzer:} „Der Kunde stellt eine Anfrage. Danach prüft ein Mitarbeiter die Angaben. 
	Wenn alles vollständig ist, wird ein Angebot erstellt. Anschließend endet der Prozess.“ \\
\end{quote}

Das System identifiziert die wesentlichen Aktivitäten, erkennt Start- und Endpunkte und leitet aus der logischen Struktur der Äußerungen eine eindeutige Sequenz ab. Die generierte JSON-Ausgabe bildet ein klares lineares Prozessmodell, das nach Verarbeitung durch den Parser und die Layoutengine als visuell kohärentes BPMN-Diagramm erscheint.

Das Beispiel zeigt, dass das System in der Lage ist, auch ohne explizite Verbindungsbegriffe (z.\,B. „dann“, „daraufhin“) die logische Abfolge eines Prozesses korrekt zu interpretieren. Dies unterstreicht die Relevanz der semantischen Interpretationsfähigkeit des LLMs, wie in Kapitel~2.2 beschrieben.

\subsubsection*{4.5.2 Beispiel 2: Entscheidungsstrukturen}

Entscheidungsprozesse stellen erhöhte Anforderungen an die Modellgenerierung, da konditionale Aussagen präzise in BPMN-Gateways übersetzt werden müssen. Ein Beispiel verdeutlicht diese Herausforderung:

\begin{quote}
	\textit{Nutzer:} „Der Antrag wird geprüft. Falls Unterlagen fehlen, informieren wir den 
	Kunden. Wenn alles vollständig ist, legen wir den Vorgang an.“ \\
\end{quote}

Das System erkennt automatisch die Existenz eines exklusiven Entscheidungspunktes und transformiert die konditionalen Aussagen in zwei alternative Sequenzflüsse. Dabei werden die Verzweigungsbedingungen in Form von Labels den jeweiligen Gateway-Ausgängen zugeordnet. Die erzeugte Repräsentation enthält bereits die korrekte logische Struktur, sodass das visuelle Modell unmittelbar interpretierbar ist.

Auffällig ist hierbei die Fähigkeit des Systems, unterschiedliche sprachliche Konstruktionen („falls…“, „wenn…“) als semantisch gleichwertig zu erkennen und konsistent einer Gateway-Struktur zuzuordnen. Dieses Verhalten entspricht genau jener LLM-gestützten Bedeutungserschließung, die in Kapitel~3.3 als zentrale Grundlage des Mappings beschrieben wurde.

\subsubsection*{4.5.3 Beispiel 3: Parallele Abläufe}

Ein weiteres Beispiel betrifft Prozesse, die parallele Aktivitäten enthalten. Solche Modelle sind in traditionellen Modellierungsszenarien häufig fehleranfällig, da die sprachliche Beschreibung paralleler Abläufe oft implizit erfolgt. In einer beispielhaften Interaktion formuliert ein Nutzer:

\begin{quote}
	\textit{Nutzer:} „Nachdem die Bestellung eingegangen ist, prüfen wir den Lagerbestand und 
	informieren gleichzeitig den Versand.“ \\
\end{quote}

Obwohl der Begriff „gleichzeitig“ lediglich linguistisch eingebettet ist, ordnet das System die Aktivitäten korrekt einem parallelen Gateway zu. Die JSON-Ausgabe enthält einen parallel\_gateway mit zwei parallelen Pfaden, was in der visuellen Darstellung zu einer klar erkennbaren Gabelung führt.

Dieses Beispiel zeigt, dass das LLM nicht lediglich strukturierende Begriffe erkennt, sondern in der Lage ist, aus impliziten sprachlichen Hinweisen eine formale Parallelstruktur abzuleiten. Damit bestätigt sich die Annahme, dass LLM-basierte Analyseverfahren in Bereichen überzeugen können, in denen regelbasierte Systeme an ihre Grenzen stoßen.

\subsubsection*{4.5.4 Diskussion der Modellqualität}

Die drei dargestellten Beispiele verdeutlichen, dass der Prototyp für eine Vielzahl typischer Modellierungssituationen konsistente und gut interpretierbare BPMN-Diagramme erzeugen kann. Dies gilt sowohl für lineare Abläufe als auch für Entscheidungs- und Parallelstrukturen. Auffällig ist vor allem die Fähigkeit des Systems, sprachliche Variabilität zu bewältigen und dennoch eine stabile Modellstruktur zu generieren. 

Die Qualität der Modelle hängt dabei maßgeblich von der Präzision der Eingaben ab, wobei das System durch seine Rückfragenmechanismen aktiv zur Klärung von Ambiguitäten beiträgt. Damit entsteht ein iterativer Interaktionsprozess, der eine kooperative Modellierung erlaubt und zugleich die Robustheit erhöht. 

Die Ergebnisse zeigen insgesamt, dass der entwickelte Ansatz geeignet ist, konzeptionelle Überlegungen aus Kapitel~3 praktisch umzusetzen und die Potenziale LLM-gestützter Modellierung im Kontext der Prozessdigitalisierung erfahrbar zu machen. Die Beispiele bilden somit eine wichtige Grundlage für die Evaluationsphase in Kapitel~5, in der die Qualität der Modelle systematisch untersucht wird.

\newpage
\makeatletter
\renewcommand\@biblabel[1]{\textbf{[#1]}} % erzwingt Anzeige der Nummer
\makeatother
\setlength{\bibhang}{1em}
\bibliographystyle{plainnat}   % oder abbrvnat, unsrtnat, dinat
\bibliography{literatur}       % Dateiname OHNE .bib


% --------------------------------------------------------------
% ABKÜRZUNGSVERZEICHNIS
% --------------------------------------------------------------
\section*{Abkürzungsverzeichnis}
\begin{acronym}[BPMN]
    % --- Grundlagen ---
    \acro{BPMN}{Business Process Model and Notation}
    \acro{LLM}{Large Language Model}
    \acro{KI}{Künstliche Intelligenz}
    \acro{AI}{Artificial Intelligence}
    \acro{ML}{Machine Learning}
    \acro{NLP}{Natural Language Processing}
    \acro{DL}{Deep Learning}
    
    % --- Prozessmanagement ---
    \acro{BPM}{Business Process Management}
    \acro{BPEL}{Business Process Execution Language}
    \acro{DMN}{Decision Model and Notation}
    \acro{UML}{Unified Modeling Language}

    % --- Systemarchitektur / Technik ---
    \acro{API}{Application Programming Interface}
    \acro{UI}{User Interface}
    \acro{UX}{User Experience}
    \acro{IDE}{Integrated Development Environment}
    \acro{JSON}{JavaScript Object Notation}
    \acro{XML}{Extensible Markup Language}

    % --- Evaluations- und Forschungsmethoden ---
    \acro{RQ}{Research Question}
    \acro{F1}{Forschungsfrage 1}
    \acro{F2}{Forschungsfrage 2}
    \acro{F3}{Forschungsfrage 3}
    \acro{KPIs}{Key Performance Indicators}

    % --- Frameworks und Werkzeuge ---
    \acro{RASA}{Rasa Open Source Framework}
    \acro{NLG}{Natural Language Generation}
    \acro{NLI}{Natural Language Interpretation}
    \acro{BERT}{Bidirectional Encoder Representations from Transformers}
    \acro{GPT}{Generative Pre-trained Transformer}
    \acro{LLAMA}{Large Language Model Meta AI}
    \acro{HUG}{Hugging Face Plattform}

    % --- Software Engineering / System Design ---
    \acro{CRUD}{Create, Read, Update, Delete}
    \acro{MVC}{Model View Controller}
    \acro{DBMS}{Database Management System}
    \acro{REST}{Representational State Transfer}
    \acro{JSON-LD}{JSON for Linked Data}

    % --- Evaluation und Tests ---
    \acro{QA}{Quality Assurance}
    \acro{PoC}{Proof of Concept}
    \acro{MVP}{Minimum Viable Product}
\end{acronym}




\end{document}



